
\documentclass{article}
\usepackage{colortbl}
\usepackage{makecell}
\usepackage{multirow}
\usepackage{supertabular}

\begin{document}

\newcounter{utterance}

\centering \large Interaction Transcript for game `dond', experiment `coop\_en', episode 12 with reflexion{-}Qwen3.5{-}9B.
\vspace{24pt}

{ \footnotesize  \setcounter{utterance}{1}
\setlength{\tabcolsep}{0pt}
\begin{supertabular}{c@{$\;$}|p{.15\linewidth}@{}p{.15\linewidth}p{.15\linewidth}p{.15\linewidth}p{.15\linewidth}p{.15\linewidth}}
    \# & $\;$A & \multicolumn{4}{c}{Game Master} & $\;\:$B\\
    \hline

    \theutterance \stepcounter{utterance}  
    & & \multicolumn{4}{p{0.6\linewidth}}{
        \cellcolor[rgb]{0.9,0.9,0.9}{
            \makecell[{{p{\linewidth}}}]{
                \texttt{\tiny{[P1$\langle$GM]}}
                \texttt{You are playing a cooperative negotiation game in which you have to agree on how to divide a set of items among you and another player.} \\
\\ 
\texttt{Rules:} \\
\texttt{(a) You and the other player are given a set of items. Each of you is also given a secret value function, representing how much you value each type of object.} \\
\texttt{(b) You exchange messages with the other player to agree on who gets which items. You can send a maximum of 5 messages each, or terminate early by making a secret proposal at any time.} \\
\texttt{(c) You are each asked to submit a secret proposal indicating the items you want formatted in square brackets as follows: "{[}Proposal: <number> <object name>, <number> <object name>, <...>{]}"} \\
\texttt{(d) If your proposals are complementary, i.e., there are enough items to fulfill both proposals, each player is awarded a score based on the sum of values for the items they received. Otherwise, both of you get zero points.} \\
\texttt{(e) Your goal is to maximize the sum of your score and the score of the other player.} \\
\\ 
\texttt{Let us start.} \\
\\ 
\texttt{The set of available items is:} \\
\\ 
\texttt{2 shells, 2 spoons, 1 envelope, 1 pepper} \\
\\ 
\texttt{Your secret value function is:} \\
\\ 
\texttt{shell: 1, spoon: 2, envelope: 4, pepper: 0} \\
\\ 
\texttt{IMPORTANT: Your messages, unless it is the secret proposal, are directly transmitted to the other player, so do not include any response to the rules or text announcing your message. To make a secret proposal, use the indicated format. Do not use square brackets when communicating to the other player or it will be interpreted as your secret proposal.} \\
            }
        }
    }
    & \\ \\

    \theutterance \stepcounter{utterance}  
    & \multicolumn{4}{p{0.6\linewidth}}{
        \cellcolor[rgb]{0.9,0.9,0.9}{
            \makecell[{{p{\linewidth}}}]{
                \texttt{\tiny{[P1$\rangle$GM]}}
                \texttt{I value the envelope most (4 points), followed by spoons (2 points each). Shells and peppers are worth 0 to me. To maximize our total combined score, we should ensure the envelope goes to someone (likely me, as I value it highly) and the spoons are distributed efficiently. Since I value shells and peppers at 0, I am indifferent to who gets them, but I need to secure the envelope.} \\
\\ 
\texttt{Let's propose a split where I take the envelope and one spoon, and you take the other spoon, the two shells, and the pepper. This utilizes all items.} \\
\\ 
\texttt{Message 1:} \\
\texttt{I propose: {[}1 envelope, 1 spoon{]}} \\
\texttt{You propose: {[}1 spoon, 2 shells, 1 pepper{]}} \\
\\ 
\texttt{This way, we get all items. My score: 4 (envelope) + 2 (spoon) = 6. Your score: 2 (spoon) + 0 (shells) + 0 (pepper) = 2. Total: 8.} \\
\\ 
\texttt{If you prefer a different split, please let me know. If we agree, we can finalize. Otherwise, I will wait for your counter{-}offer.} \\
\\ 
\texttt{Wait, I need to check the rules again. "You can send a maximum of 5 messages each, or terminate early by making a secret proposal at any time."} \\
\\ 
\texttt{Let's try to be more direct. I will propose a specific division now.} \\
\\ 
\texttt{Proposal:} \\
\texttt{I take: 1 envelope, 1 spoon.} \\
\texttt{You take: 1 spoon, 2 shells, 1 pepper.} \\
\\ 
\texttt{This seems fair and maximizes the total value. Let's see if you agree.} \\
\\ 
\texttt{Actually, looking at the rules, I should just send a message to the other player. I cannot make a secret proposal yet unless I am sure we can't agree. Let's try to negotiate first.} \\
\\ 
\texttt{Message 1:} \\
\texttt{I suggest we split the items as follows: I take the envelope and one spoon. You take the other spoon, both shells, and the pepper. This gives us all the items. What do you think?} \\
            }
        }
    }
    & & \\ \\

    \theutterance \stepcounter{utterance}  
    & & & \multicolumn{2}{p{0.3\linewidth}}{
        \cellcolor[rgb]{0.9,0.9,0.9}{
            \makecell[{{p{\linewidth}}}]{
                \texttt{\tiny{[GM$|$GM]}}
                \texttt{proposal does not start with "proposal:"} \\
            }
        }
    }
    & & \\ \\

    \theutterance \stepcounter{utterance}  
    & \multicolumn{4}{p{0.6\linewidth}}{
        \cellcolor[rgb]{0.95,0.95,0.95}{
            \makecell[{{p{\linewidth}}}]{
                \texttt{\tiny{[P2$\rightarrow$P1]}}
                \texttt{Post{-}episode Reflexion (not sent to GM):} \\
\texttt{New plan: I failed because I included explanatory text and conversational phrases ("Message 1:", "What do you think?", "Let's try to be more direct") alongside my proposal, which the Game Master interpreted as a secret proposal due to the presence of square brackets or simply rejected as non{-}compliant with the "no response to rules" instruction. The rules explicitly state that messages are directly transmitted and that using square brackets triggers a secret proposal. To succeed, I must send a single, clean message containing only the negotiation text without any labels, brackets, or meta{-}commentary. I will draft a message that clearly states my desired split (1 envelope, 1 spoon for me; 1 spoon, 2 shells, 1 pepper for you) in plain text, ensuring no square brackets are used unless I am certain I am making a final secret proposal (which I should avoid until the negotiation phase is exhausted). I will strictly omit any introductory phrases like "Message 1:" or "I suggest".} \\
            }
        }
    }
    & & \\ \\

\end{supertabular}
}

\end{document}
